% ===================================================================
%  도표 제 4 호 — 한창때와 늘그막.
%
%  Traduction, faite en 2026, en coreen d'aujourd'hui, sur l'ido de
%  texto/io. Regles de la colonne : voir 10-dopyo-01.tex. Deux
%  scenes, deux numerotations : les cles portent c1 et c2.
%
%  LA PARENTE EST LE VRAI SUJET DE CE TABLEAU, ET LE COREEN Y EST
%  PLUS PRECIS QUE L'IDO. La noce reunit deux familles, et l'ido s'en
%  tire avec des composes generiques — bopatrulo, bomatro, bofilio,
%  bofiliino, bofrato, bofratino — qui laissent au lecteur le soin de
%  deviner de quel cote. Le coreen ne le permet pas : chaque lien a
%  son mot, et il faut trancher a chaque fois.
%
%    « Elua bopatrulo (20) » est le pere du MARI, donc 시아버지 ;
%    « la bomatro (21) di Ioannes » est la mere de la FEMME, donc
%    장모 ; la bru est 며느리, la belle-sœur du bloc c1-13-1 est
%    시누이 — la sœur du mari — et le beau-frere du bloc c1-15-1 est
%    동서, le mot qui designe le mari de la sœur de sa femme et rien
%    d'autre. Cinq mots la ou l'ido en a cinq aussi, mais qui ne se
%    recouvrent pas : ce sont les liens du fac-simile qui ont
%    tranche, lu phrase par phrase.
%
%  ET 들러리 EST EXACTEMENT LE GARCON ET LA DEMOISELLE D'HONNEUR
%  (25, 26) : un seul mot coreen pour les deux, ce que le francais
%  fait aussi et que l'ido ne fait pas.
%
%  DIX RETOURNEMENTS SUR DIX-NEUF PAIRES. Le tableau tient le milieu
%  entre le defile du tableau 3 (dix-sept sur vingt-cinq) et
%  l'enumeration du tableau 2 (un sur sept) : la premiere scene
%  enumere les convives, la seconde les decrit par attributs — le
%  manteau AVEC son col de fourrure, le coussin POSE sur le tabouret,
%  la robe de chambre AVEC ses gros boutons, l'ordonnance QUE le
%  medecin vient d'ecrire.
%
%  DEUX CHOSES QUE LE COREEN NE NOMME PAS D'UN MOT, et je prefere le
%  dire : le « kastelo » (1) du bloc c1-03-1 — un manoir de campagne
%  francais, que 성 rendrait par une forteresse et que la colonne
%  ecrit 저택, la grande maison — et le « verando » (8), qui s'ecrit
%  베란다, l'emprunt courant. Le coreen a bien 툇마루, la galerie de
%  la maison coreenne, mais c'est un autre objet, et l'ecrire ici
%  deplacerait la maison entiere.
% ===================================================================

%%K t04-apar-1 apar
\VUsaut{4.0mm}
\VUtitre{15.0pt}{60}{도표 제 4 호}
\VUsaut{1.6mm}
\VUtitre{11.4pt}{0}{한창때와 늘그막.}
\VUsaut{3.2mm}

%%K t04-tit-1 sub
\VUcentre{11.4pt}{0}{\textit{첫째 마당.}}
\VUsaut{1.6mm}
\VUtitre{11.4pt}{0}{혼인 잔치.}
\VUsaut{2.6mm}

%%K t04-tit-2 sub
\VUpk{10.2pt}{40}{가을.}
\VUsaut{3.0mm}

%%K t04-c1-01-1 p
1. --- 젊은이들이 자랐다. 그 가운데 하나인 요안네스가 혼인한다. 우리는
\VUgras{혼인 잔치}에 함께한다. 그 잔치는 신랑 신부의 두 집안을 한 상에 모은다.

%%K t04-c1-02-1 p
2. --- 지금은 \VUgras{가을}, \VUgras{구월}이다. \VUgras{시월}과 \VUgras{십일월}의
찬바람은 아직 들판에서 푸른 잎을 벗겨 가지 않았다. 저녁 공기는 미지근하고 향기롭다.
그래서 저녁상은 뜰에 면한 \VUgras{베란다} \textsuperscript{(8)} 아래에 차려진다.

%%K t04-c1-03-1 p
3. --- 멀리 \VUgras{언덕} \textsuperscript{(3)} 위에 요안네스 부모님의 집이 있다. 푸른
숲에 가려진 맵시 있는 \VUgras{저택} \textsuperscript{(1)}이다. 좋은 포도주를 내는 고운
포도밭이 그 언덕의 비탈을 덮고 있다. 포도 농부가 그것을 가꾸고 돌본다.

%%K t04-c1-04-1 p
4. --- 포도밭 위로 \VUgras{자고새} \textsuperscript{(2)} 떼가 날아간다 —
\VUgras{사냥터지기} \textsuperscript{(5)}가 다가오자 놀란 것이다. 그 사람은 겨드랑이에
\VUgras{엽총}을 끼고 어깨에 \VUgras{사냥 자루}를 메고서 \VUgras{덤불}
\textsuperscript{(4)}을 뒤진다 — 제 \VUgras{개} \textsuperscript{(6)}와 함께. 그는 땅을
지키고 밀렵꾼들이 사냥감을 없애지 못하게 막는다.

%%K t04-c1-05-1 p
5. --- \VUgras{베란다} 바깥으로 \VUgras{포도 시렁}이 뻗어 오르는데, 거기에 굵은
\VUgras{포도송이} \textsuperscript{(7)}가 매달려 있다.

%%K t04-c1-06-1 p
6. --- \VUgras{상} \textsuperscript{(16)}을 꾸미는 고운 가을꽃은 \VUgras{국화}
\textsuperscript{(9)}다. 신부의 \VUgras{아버지} \textsuperscript{(10)}는 그 가운데 한 송이를
자기 연미복 단춧구멍에 꽂았다.

%%K t04-c1-07-1 p
7. --- 식사가 끝나 간다. \VUgras{하인} \textsuperscript{(11)}이 후식 접시를 나르기
시작하고, 곧 더는 쓰지 않을 상차림을 치울 것이다 — \VUgras{숟가락}
\textsuperscript{(14)}, \VUgras{포크} \textsuperscript{(12)}, \VUgras{나이프}
\textsuperscript{(13)}, \VUgras{큰 접시} \textsuperscript{(15)}를.

%%K t04-tit-3 sub
\VUpk{10.2pt}{40}{친척들.}
\VUsaut{3.0mm}

%%K t04-c1-08-1 p
8. --- 모두가 즐거워 보인다. 신랑의 \VUgras{어머니} \textsuperscript{(17)}가 제 아이들을
바라보는 그 미소가 그분의 행복을 말해 준다.

%%K t04-c1-09-1 p
9. --- 이 혼인은 고을의 넉넉한 두 집안을 잇는다. \VUgras{신랑} \textsuperscript{(18)}
요안네스는 큰 지주의 \VUgras{아들}이고, 이제 솜씨 좋은 공장주의 \VUgras{사위}가 된다.

%%K t04-c1-10-1 p
10. --- \VUgras{부부}는 젊지만 오래전부터 \VUgras{약혼한} 사이였다.
\VUgras{새색시} \textsuperscript{(19)} 마르타는 귤꽃으로 꾸민 \VUgras{흰 드레스} 차림이
아주 곱다.

%%K t04-c1-11-1 p
11. --- 그 여자의 \VUgras{시아버지} \textsuperscript{(20)}는 그토록 상냥한
\VUgras{며느리}를 얻어 아주 행복하다. 요안네스의 \VUgras{장모} \textsuperscript{(21)}는
제 \VUgras{딸}이 그토록 고운 것을 보고 미소 짓는다.

%%K t04-c1-12-1 p
12. --- 상 끝에는 나머지 \VUgras{손님들} \textsuperscript{(22)}이 앉아 있다 —
\VUgras{친척}, \VUgras{벗}, \VUgras{사촌 형제}, \VUgras{사촌 자매}다.
\VUgras{잔치 책임자} \textsuperscript{(23)}가 겨드랑이에 행주를 끼고 부지런히 시중을 든다.

%%K t04-tit-4 sub
\VUpk{10.2pt}{40}{벗들.}
\VUsaut{3.0mm}

%%K t04-c1-13-1 p
13. --- 상 오른쪽에는 \VUgras{아가씨} \textsuperscript{(24)} 하나가 있다. 새색시의
\VUgras{벗}이다. 그 다음이 그 아가씨의 \VUgras{오라비}인데, 그가 새색시의
\VUgras{들러리} \textsuperscript{(25)}다. 그는 제 \VUgras{신부 들러리} \textsuperscript{(26)}와
이야기를 나눈다. 그 여자는 신랑의 누이이고, 이 혼인으로 마르타의 \VUgras{시누이}가
된다. 그 여자는 고운 젊은이다. \VUgras{패물}과 \VUgras{진주 목걸이}, 금
\VUgras{팔찌}를 지녔다.

%%K t04-c1-14-1 p
14. --- 그들 곁에서 일어서서 잔을 드는 어른은 집안의 \VUgras{벗} \textsuperscript{(27)}이다.
그는 \VUgras{신랑 쪽 증인} 가운데 하나다. 그는 \VUgras{축배를 들어} 젊은 부부의
\VUgras{건강}을 위해 마시고 오랜 행복을 빈다. 그는 예복 차림에 연미복을 입고
\VUgras{에나멜 구두} \textsuperscript{(28)}를 신었다. 그는 \VUgras{할머니}
\textsuperscript{(29)}를 모시고 온 사람인데, 그 할머니는 \VUgras{홀로 되신} 분이다.

%%K t04-c1-15-1 p
15. --- \VUgras{연미복} \textsuperscript{(31)} 차림에 검고 가는 콧수염을 기른 젊은이는
요안네스의 \VUgras{동서} \textsuperscript{(30)}다. 그는 이름난 기술자다. 그는 아주 맵시
있는 \VUgras{젊은 부인} \textsuperscript{(33)} 곁에 앉아 있는데, 그 부인은 마르타의
\VUgras{언니}이고, 사촌 오빠에게 팔을 내밀고 다가와 꽃을 건네며
\VUgras{안녕히 주무시라고} \VUgras{빌어} 주는 귀여운 계집아이의 어머니다. 그 부인은
아주 고운 \VUgras{귀고리}와 맵시 있는 \VUgras{머리 장식}을 하고 있다.

%%K t04-c1-16-1 p
16. --- 어린 사촌 오빠는 제 \VUgras{덧옷} \textsuperscript{(36)}과 불룩한
\VUgras{반바지} \textsuperscript{(37)} 때문에 아주 어색해 보이는데, 참 딱하다. 그는
\VUgras{고아} \textsuperscript{(35)}다. 아주 어려서 아버지와 어머니를 잃었다. 그의
아저씨가 \VUgras{후견인} \textsuperscript{(38)}이 되어, 그 딱한 아이를 기르려고 홀몸으로
남았다. 마음씨 좋은 어른이다. 머리가 좀 벗어졌고 눈이 아주 나빠서 \VUgras{코안경}을
쓴다.

%%K t04-tit-5 sub
\VUsaut{2.4mm}
\VUpk{10.2pt}{40}{후식.}
\VUsaut{3.0mm}

%%K t04-c1-17-1 p
17. --- 잔치 손님들 뒤로, \VUgras{상보}를 덮은 상 위에 \VUgras{후식}
\textsuperscript{(39)}이 차려져 있다. 큰 굽접시에 과일이 있다 — \VUgras{복숭아}
\textsuperscript{(40)}, \VUgras{배} \textsuperscript{(41)}, \VUgras{사과}
\textsuperscript{(42)}, \VUgras{살구} \textsuperscript{(43)}, \VUgras{포도}
\textsuperscript{(44)}다. 그 곁에는 \VUgras{파인애플} \textsuperscript{(45)}, 고운
\VUgras{케이크} \textsuperscript{(46)}, \VUgras{술병} \textsuperscript{(48)},
\VUgras{샴페인} \textsuperscript{(49)}, 그리고 깨끗한 \VUgras{접시}
\textsuperscript{(47)} 무더기가 있다.

%%K t04-c1-18-1 p
18. --- \VUgras{술창고지기} \textsuperscript{(50)}가 오래된 포도주 \VUgras{병}
\textsuperscript{(52)}을 딴다 — \VUgras{코르크 따개} \textsuperscript{(51)}로.
\VUgras{코르크 마개} \textsuperscript{(53)}는 몹시 뽑기 어려워 보인다. 그 사람은
겨드랑이에 \VUgras{냅킨} \textsuperscript{(54)}을 끼고 있다.

%%K t04-c1-19-1 p
19. --- 저녁을 마치면 어른들은 담배 피우는 방으로 가고, 부인들은 무도회 채비를 하러 갈
것이다.

%%K t04-tit-6 sub
\VUsaut{5.0mm}
\VUcentre{11.4pt}{0}{\textit{둘째 마당.}}
\VUsaut{1.6mm}
\VUpk{11.4pt}{40}{할아버지 생신.}
\VUsaut{2.6mm}

%%K t04-tit-7 sub
\VUpk{10.2pt}{40}{문안.}
\VUsaut{3.0mm}

%%K t04-c2-01-1 p
1. --- 열 해가 지났다. 요안네스의 부모님은 늙으셔서 손주 셋을 몹시 사랑하신다.
\VUgras{손자} \textsuperscript{(2)} 둘과 \VUgras{손녀} \textsuperscript{(3)} 하나다. 오늘이
\VUgras{할아버지 생신}이라 아이들이 생신을 축하드리러 왔다.

%%K t04-c2-02-1 p
2. --- 어린 페트루스는 아직 아주 어리다. 세 살이다. 그는 \VUgras{고양이}
\textsuperscript{(1)}가 다가오는 것을 본다. 그는 그 고운 비단 같은 \VUgras{털}과 긴
\VUgras{꼬리}를 쓰다듬고 싶어 한다. 그 얌전한 \VUgras{발} 밑에 제 몸을 할퀼지도 모를
날카롭고 짓궂은 발톱이 숨어 있다는 것은 생각하지 않는다.

%%K t04-c2-03-1 p
3. --- 그에게 손을 내미는 누이 가브리엘라는 훨씬 의젓하다. 마르켈루스는 큰
\VUgras{외투} \textsuperscript{(4)}에 가죽 \VUgras{허리띠} \textsuperscript{(5)} 차림이라
제법 어른스러워 보인다. 짧은 반바지 밑으로 \VUgras{양말} \textsuperscript{(6)}이 보이기는
하지만.

%%K t04-c2-04-1 p
4. --- 그들의 아버지는 두꺼운 겨울 \VUgras{외투} \textsuperscript{(7)}를 입으셨다 —
\VUgras{털 깃} \textsuperscript{(8)}이 달린 것이다 —, 그리고 \VUgras{주머니}
\textsuperscript{(9)}에는 목도리를 넣으셨다. 그 부인은 \VUgras{깃털}
\textsuperscript{(10)}로 꾸민 펠트 모자를 쓰고, \VUgras{재킷} \textsuperscript{(11)},
\VUgras{목도리 털} \textsuperscript{(12)}, \VUgras{토시} \textsuperscript{(13)}를 갖추었다.

%%K t04-c2-05-1 p
5. --- 아주 춥다. 겨울이 한창이기 때문이다. \VUgras{눈} \textsuperscript{(19)}이 온종일
내려서 집집의 지붕이 온통 하얗다. \VUgras{밤}과 함께 추위는 더욱 매워진다. 하늘에서는
\VUgras{달} \textsuperscript{(17)}과 \VUgras{별} \textsuperscript{(18)}이 빛나
\VUgras{어둠}을 꿰뚫는다.

%%K t04-tit-8 sub
\VUsaut{4.2mm}
\VUpk{10.2pt}{40}{병.}
\VUsaut{3.0mm}

%%K t04-c2-06-1 p
6. --- 딱한 \VUgras{앞 못 보는 이} \textsuperscript{(20)}가 \VUgras{약방}
\textsuperscript{(16)} 앞 광장을 지나간다 — 제 \VUgras{푸들} \textsuperscript{(21)}이
이끄는 대로. 그는 참으로 불행하다. 유스티니아라는 \VUgras{하녀} \textsuperscript{(14)}는 부엌에서
아주 뜨거운 \VUgras{약차}가 담긴 \VUgras{사발} \textsuperscript{(15)}을 가져온다. 설탕을
넉넉히 넣은 것이다.

%%K t04-c2-07-1 p
7. --- \VUgras{할머니} \textsuperscript{(23)}가 상 위의 \VUgras{손잡이 안경}
\textsuperscript{(26)}을 찾으려 하신다. \VUgras{처방} \textsuperscript{(27)}을 읽으시려는
것인데, 그것은 \VUgras{의사} \textsuperscript{(25)}가 방금 써 준 것이다. 그분은 제
\VUgras{지팡이} \textsuperscript{(24)}에 몸을 기대고 안락의자에서 일어나신다.

%%K t04-c2-08-1 p
8. --- 그분은 자주 \VUgras{잔병}, \VUgras{어지럼}, \VUgras{고뿔}, \VUgras{치통}으로
앓으신다. 다리는 약하고 눈도 이제 그리 좋지 않다. 그런데도 그분은 늘 \VUgras{물약}
\textsuperscript{(28)}을 지어 주려 드는 늙은 \VUgras{의사}를 즐겨 놀리신다. 오늘 그분은
젊은 식구들이 둘러앉아 아주 즐거우시다.

%%K t04-c2-09-1 p
9. --- 잘 닫힌 방 안은 아늑하다. 대리석 \VUgras{벽난로} \textsuperscript{(29)}에는
\VUgras{활활 타는 불} \textsuperscript{(30)}이 피어오르고, 방금 켠 \VUgras{등}
\textsuperscript{(31)}의 부드러운 \VUgras{불빛} 속에서 사람들은 행복을 느낀다.

%%K t04-tit-9 sub
\VUsaut{4.2mm}
\VUpk{10.2pt}{40}{할아버지.}
\VUsaut{3.0mm}

%%K t04-c2-10-1 p
10. --- \VUgras{할아버지} \textsuperscript{(33)}조차 잊으셨다. 당신이 앓고 계셨다는 것,
\VUgras{신경통} 때문에 \VUgras{안락의자} \textsuperscript{(34)}에 붙박여 계신다는 것,
어제만 해도 \VUgras{열과 실신}이 있었다는 것을. 지난겨울에 \VUgras{폐렴}을 앓으셨고
쉰 목소리로 힘든 \VUgras{기침}을 하시느라 몹시 괴로우셨다는 것도 이제 기억하지 않으신다.

%%K t04-c2-10-2 p
그런데도 병은 아주 어렵게 나았다. 이제 좀 나으시다. 다만 \VUgras{방석}
\textsuperscript{(38)} 위에 올려놓으신 다리가 여전히 아프시다 — 그 방석은 작은
\VUgras{발판} \textsuperscript{(39)} 위에 놓여 있다.

%%K t04-c2-11-1 p
11. --- 따뜻한 \VUgras{실내복} \textsuperscript{(35)}에 몸을 잘 감싸실 때 — 자개로 된
굵은 \VUgras{단추} \textsuperscript{(36)}가 달린 옷이다 —, 그리고 곁에서 신문을 읽어 주는
어린 \VUgras{고아 소녀} \textsuperscript{(40)}가 있을 때, 그분은 투정하지 않으신다. 그
아이는 흰 \VUgras{머리쓰개}를 쓰고 푸른 \VUgras{옷} \textsuperscript{(42)}과
\VUgras{어깨망토} \textsuperscript{(41)}를 걸치고 있다.

%%K t04-c2-12-1 p
12. --- 해마다 \VUgras{달력} \textsuperscript{(43)}이 \VUgras{십이월} 초하루
\VUgras{날짜}를 다시 가리킬 때면, 그분은 \VUgras{손주들}을 애타게 기다리신다. 그
아이들이 그 중요한 \VUgras{날짜}를 잊지 않으리라는 것을 아시기 때문이다.

%%K t04-c2-13-1 p
13. --- 그분은 아주 오래전, 몸소 황제의 이름난 \VUgras{원수}께 생신 인사를 드리러 가던
때를 뭉클하게 기억하신다. 그분의 \VUgras{초상} \textsuperscript{(44)}이 벽에 걸려 있고,
그분이 당신 아이들의 \VUgras{증조할아버지}시다.

\VUsaut{6.0mm}
\VUfilet{22mm}
